\documentclass[12pt,letterpaper]{article}

\usepackage[utf8]{inputenc}
\usepackage[T1]{fontenc}
\usepackage[spanish,es-nodecimaldot]{babel}
\usepackage{geometry}
\usepackage{amsmath}
\usepackage{booktabs}
\usepackage{tabularx}
\usepackage{enumitem}
\usepackage[hidelinks]{hyperref}

\geometry{margin=2.5cm}
\setlength{\parindent}{0pt}
\setlength{\parskip}{0.7em}
\renewcommand{\arraystretch}{1.25}

\title{Documentación del proyecto\\[0.3em]
\large Guía interactiva para la configuración de routers Digi}
\author{Nombre del autor o equipo}
\date{\today}

\begin{document}

\maketitle

\begin{abstract}
Este proyecto consiste en una aplicación web estática orientada a la capacitación
técnica y a la consulta rápida de procedimientos de configuración para routers
Digi. La solución presenta páginas HTML independientes con instrucciones visuales
y recorridos paso a paso para configurar Syslog, autenticación, grupos de usuarios,
RADIUS y TACACS+. Su propósito es simplificar el aprendizaje, reducir el tiempo
invertido en localizar opciones dentro de la interfaz del equipo y disminuir los
errores asociados con la interpretación de documentación extensa.
\end{abstract}

\section{Tipo de proyecto}

El proyecto se clasifica como una \textbf{aplicación web estática de capacitación
técnica y apoyo operativo}. También puede describirse como una guía interactiva o
un centro de aprendizaje de autoservicio para la configuración de routers Digi.

Es una aplicación estática porque utiliza páginas HTML y recursos gráficos locales,
sin requerir una base de datos, un servidor de aplicaciones ni un proceso de
instalación. El usuario solamente necesita abrir el archivo \texttt{index.html} en
un navegador web moderno. Debido a este diseño, el contenido puede utilizarse de
forma local, compartirse con otros técnicos o publicarse en un servidor web
sencillo.

\section{Problema que atiende}

La configuración de servicios de autenticación, autorización y registro en un
router puede requerir consultar manuales extensos, buscar la ubicación exacta de
cada opción y alternar repetidamente entre la documentación y la interfaz del
dispositivo. Esto incrementa el tiempo de aprendizaje y puede provocar omisiones,
especialmente cuando el técnico no realiza estas tareas de manera frecuente.

El proyecto concentra los procedimientos principales en un mismo sitio y los
convierte en recorridos visuales breves. De esta manera, el usuario puede identificar
qué opción debe modificar, qué valor debe introducir y en qué orden debe completar
el procedimiento.

\section{Objetivo general}

Proporcionar una herramienta visual, accesible y fácil de consultar que apoye la
capacitación y la preparación de técnicos para configurar servicios esenciales en
routers Digi.

\subsection{Objetivos específicos}

\begin{itemize}[leftmargin=*]
    \item Centralizar en una sola interfaz los procedimientos de configuración más
    importantes.
    \item Mostrar cada procedimiento mediante instrucciones y pasos visuales.
    \item Facilitar el aprendizaje autónomo sin depender constantemente de un
    instructor.
    \item Reducir el tiempo dedicado a buscar opciones en manuales o en la interfaz
    del router.
    \item Servir como referencia rápida antes de trabajar con un dispositivo real.
    \item Promover configuraciones más consistentes al presentar siempre la misma
    secuencia de trabajo.
\end{itemize}

\section{Alcance funcional}

La página principal funciona como un panel de navegación hacia cinco módulos de
capacitación. En conjunto, los módulos contienen 21 pasos guiados.

\begin{table}[htbp]
\centering
\small
\begin{tabularx}{\textwidth}{@{}p{2.7cm}p{1.5cm}X@{}}
\toprule
\textbf{Módulo} & \textbf{Pasos} & \textbf{Función} \\
\midrule
Syslog & 3 & Explica cómo habilitar el envío de eventos, indicar el servidor remoto y configurar el puerto y el protocolo de comunicación. \\
Autenticación & 4 & Presenta cómo agregar y ordenar métodos de autenticación mediante RADIUS, TACACS+ y usuarios locales. \\
Grupos de usuarios & 3 & Muestra cómo crear un grupo y asignar permisos y niveles de acceso de forma uniforme. \\
RADIUS & 5 & Guía la definición del servidor, puerto, secreto compartido y opciones de comportamiento y diagnóstico. \\
TACACS+ & 6 & Explica la configuración del servidor, puerto, secreto, atributos de grupo, servicio y modo autoritativo. \\
\bottomrule
\end{tabularx}
\caption{Módulos incluidos en la aplicación.}
\label{tab:modulos}
\end{table}

Cada módulo combina una imagen de la interfaz del router con una zona resaltada y
un panel de instrucciones. Los pasos avanzan automáticamente, pero el usuario puede
pausar la reproducción, reanudarla o desplazarse hacia adelante y hacia atrás. Esta
interacción permite estudiar el procedimiento al ritmo de cada persona.

\section{Funcionamiento de la solución}

El flujo normal de uso es el siguiente:

\begin{enumerate}[leftmargin=*]
    \item El usuario abre \texttt{index.html} en un navegador.
    \item Selecciona el módulo correspondiente a la tarea que desea aprender o
    consultar.
    \item La página muestra instrucciones, valores de referencia y una imagen de la
    interfaz del router.
    \item Un indicador visual resalta la sección que requiere atención en cada paso.
    \item El usuario pausa, repite o adelanta el recorrido según sea necesario.
    \item Con la secuencia comprendida y los datos del entorno confirmados, el
    técnico puede realizar la configuración en el equipo real.
\end{enumerate}

La aplicación es exclusivamente instructiva. No se conecta al router, no almacena
credenciales, no envía comandos y no modifica la configuración de ningún equipo.
Esta separación permite utilizarla como material de práctica sin poner en riesgo un
dispositivo de producción.

\section{Utilidad y beneficios}

\subsection{Capacitación más clara}

La combinación de instrucciones breves, imágenes y resaltado visual facilita la
relación entre la explicación técnica y el lugar exacto donde debe realizarse una
acción. Esto resulta especialmente útil para personal nuevo o para técnicos que
necesitan repasar un procedimiento que utilizan con poca frecuencia.

\subsection{Referencia rápida durante la preparación}

Los módulos permiten revisar un procedimiento antes de iniciar una ventana de
mantenimiento. El técnico puede confirmar la secuencia y preparar previamente
datos como direcciones IP, puertos, secretos compartidos, atributos y permisos.

\subsection{Estandarización}

Al presentar una secuencia uniforme, la herramienta ayuda a que diferentes personas
sigan el mismo proceso de capacitación. Esto reduce la dependencia de explicaciones
verbales y facilita la transferencia de conocimiento dentro del equipo.

\subsection{Portabilidad y bajo costo de operación}

La solución no necesita instalación ni infraestructura especializada. Puede abrirse
sin conexión a Internet y copiarse a otro equipo conservando su funcionamiento. Al
no existir un componente de servidor o una base de datos, el mantenimiento técnico
es reducido.

\section{Estimación del ahorro de tiempo}

El ahorro se concentra en la etapa de \textbf{búsqueda, orientación y aprendizaje};
no elimina el tiempo necesario para introducir valores, validar conectividad o
probar la configuración en el router real.

Como escenario de referencia, localizar y comprender un procedimiento mediante
documentación convencional puede tomar entre 10 y 20 minutos. Revisar el mismo tema
con un módulo visual concentrado puede tomar aproximadamente entre 3 y 7 minutos.
En un escenario medio de 15 minutos frente a 5 minutos, la herramienta ahorra unos
10 minutos por consulta, equivalentes a una reducción aproximada del 67\% en el
tiempo de orientación.

\begin{table}[htbp]
\centering
\small
\begin{tabularx}{\textwidth}{@{}Xccc@{}}
\toprule
\textbf{Actividad} & \textbf{Método tradicional} & \textbf{Guía interactiva} & \textbf{Ahorro estimado} \\
\midrule
Consulta de un tema & 10--20 min & 3--7 min & 5--10 min \\
Repaso de los cinco módulos & 50--100 min & 15--35 min & 35--65 min \\
\bottomrule
\end{tabularx}
\caption{Estimación del tiempo dedicado a localizar y comprender los procedimientos.}
\label{tab:ahorro}
\end{table}

Estas cifras son una estimación operativa y no el resultado de una prueba formal de
usuarios. Para obtener una métrica propia de la organización se recomienda registrar
el tiempo que varios técnicos necesitan para completar la misma tarea con y sin la
guía, y calcular el ahorro con la ecuación:

\[
\text{Ahorro porcentual} =
\frac{T_{\text{tradicional}} - T_{\text{guía}}}
     {T_{\text{tradicional}}} \times 100.
\]

Además del tiempo directo, la herramienta puede reducir interrupciones a personal
experimentado, repeticiones de capacitación y errores ocasionados por saltarse un
paso. Estos beneficios indirectos deben evaluarse por separado según el volumen de
usuarios y la frecuencia de las configuraciones.

\section{Usuarios previstos y casos de uso}

La solución está dirigida principalmente a técnicos de soporte, administradores de
red, personal en capacitación, instructores y estudiantes de redes. Sus casos de uso
principales son:

\begin{itemize}[leftmargin=*]
    \item inducción de nuevos integrantes del equipo;
    \item repaso previo a una configuración o demostración;
    \item apoyo visual durante una sesión de capacitación;
    \item consulta rápida de pasos y parámetros;
    \item material de estudio disponible sin conexión.
\end{itemize}

\section{Limitaciones y consideraciones de seguridad}

La guía muestra procedimientos y valores de ejemplo, pero no sustituye la
documentación oficial del modelo y la versión de firmware utilizados. Antes de
aplicar una configuración en producción se deben verificar las direcciones IP, los
puertos, los secretos compartidos, los permisos, las políticas de seguridad y las
opciones disponibles en el equipo específico.

Los secretos o credenciales reales no deben incorporarse en las páginas HTML ni en
capturas destinadas a capacitación. También es necesario validar la configuración
en un entorno controlado y mantener un mecanismo de recuperación antes de efectuar
cambios que puedan afectar el acceso administrativo.

\section{Conclusión}

La guía interactiva para routers Digi transforma procedimientos técnicos en módulos
visuales breves y accesibles. Su principal valor es centralizar el conocimiento,
acelerar la orientación del usuario y ofrecer una experiencia de capacitación
uniforme sin requerir infraestructura adicional. Aunque la estimación de ahorro debe
validarse con mediciones internas, el proyecto tiene el potencial de disminuir de
forma considerable el tiempo de consulta y de mejorar la preparación de los
técnicos antes de intervenir un equipo real.

\end{document}
